\SourceLeaf{389}%
[Differentiirt man nemlich die beiden Seiten der Identität
\[
\begin{aligned}
mmnn&=A'A'+B'B'+C'C'+A''A''+B''B''+C''C''\\
&\quad+2A(B'C''-C'B'')+2B(C'A''-A'C'')+2C(A'B''-B'A'')
\end{aligned}
\]
nach $t$ und ordnet rechts nach $A^{(2)},B^{(2)},C^{(2)}$; $A^{(3)},B^{(3)},C^{(3)}$, so werden vermöge der Gleichungen 12 und 13 die Coefficienten der Reihe nach gleich $2na,2nb,2nc$; $2n\alpha,2n\beta,2n\gamma$. Entsprechend ergibt sich durch Differentiation nach $u$:]
\[
mmn''+mnm''=aA^{(3)}+bB^{(3)}+cC^{(3)}+\alpha A^{(4)}+\beta B^{(4)}+\gamma C^{(4)},
\]
also [wird]
\begin{equation}\tag*{$(\odot)$}
\left\{\begin{aligned}
mnm'&=\alpha A^{(3)}+\beta B^{(3)}+\gamma C^{(3)}-aA^{(4)}-bB^{(4)}-cC^{(4)}\\
mnm''&=-\alpha A^{(2)}-\beta B^{(2)}-\gamma C^{(2)}+aA^{(3)}+bB^{(3)}+cC^{(3)}.
\end{aligned}\right.
\end{equation}
\par\indent Noch verdient bemerkt zu werden [dass aus der Differentiation von 12 und 13 unter Benutzung der Gleichungen $\FourRingCross$ hervorgeht:]
\[
\begin{aligned}
na'+an'&=A^{(2)}-A\Delta+CB^{(3)}-BC^{(3)}\\
n\alpha''+\alpha n''&=A^{(4)}-A\Delta+BC^{(3)}-CB^{(3)},
\end{aligned}
\]
also
\[
n'a+n''\alpha+n(\alpha''+a')=A^{(2)}+A^{(4)}-2A\Delta,
\]
oder
\[
\begin{aligned}
A^{(2)}+A^{(4)}&=n'a+n''\alpha+(2\Delta-mmnn)A\\
B^{(2)}+B^{(4)}&=n'b+n''\beta+(2\Delta-mmnn)B\\
C^{(2)}+C^{(4)}&=n'c+n''\gamma+(2\Delta-mmnn)C.
\end{aligned}
\]
\Needspace{10\baselineskip}\begin{center}[4.]\end{center}
\noindent\hspace*{1.2em} Für die [auf die Ebene] abwickelungsfähige Fläche findet man aus anderweitigen Betrachtungen, dass, wenn $Q,Q'$ willkürliche Functionen von $q$ bedeuten und aus der Auflösung [der] Gleichung
\[
qt+Qu=1
\]
